% Innovation flowchart; yellow denotes novel data and wedge boxes are neutral.
\begin{center}
\begin{tikzpicture}[
    processbox/.style={
        draw,
        rounded corners,
        align=center,
        minimum width=2.3cm,
        minimum height=0.85cm
    },
    measured/.style={processbox, fill=yellow!35},
    wedgebox/.style={
        draw,
        align=center,
        minimum width=2.05cm,
        minimum height=0.75cm,
        font=\scriptsize
    },
    processarrow/.style={-{Latex[length=2.5mm]}, thick},
    wedgearrow/.style={-{Latex[length=2mm]}, thin}
]

\node[processbox] (cost) {Financial Cost\\of Capital};
\node[measured, right=1.15cm of cost]
    (perceived) {Perceived Cost\\of Capital};
\node[measured, right=1.15cm of perceived]
    (discount) {Discount\\Rate};
\node[processbox, right=1.15cm of discount]
    (investment) {Investment};

\coordinate (mid1) at ($(cost.east)!0.5!(perceived.west)$);
\coordinate (mid2) at ($(perceived.east)!0.5!(discount.west)$);

\draw[processarrow] (cost) -- (perceived);
\draw[processarrow] (perceived) -- (discount);
\draw[processarrow] (discount) -- (investment);

\node[wedgebox, below=0.9cm of mid1] (costwedge)
    {Cost of Capital\\Wedge};
\node[wedgebox, below=0.9cm of mid2] (discountwedge)
    {Discount Rate\\Wedge};

\draw[wedgearrow] (costwedge.north) -- (mid1);
\draw[wedgearrow] (discountwedge.north) -- (mid2);

\node[
    draw,
    fill=yellow!35,
    minimum width=0.45cm,
    minimum height=0.35cm,
    below=0.55cm of discountwedge,
    xshift=-0.85cm
] (legendbox) {};
\node[
    right=0.15cm of legendbox,
    anchor=west,
    font=\small
] {Novel data in paper};

\end{tikzpicture}
\end{center}
